% GENERATED FILE. Edit canonical lessons, then run npm run generate:books.
% canonical-source-hash: fnv1a64:71d6c7ead483052c
% canonical-lessons: TE-C390-vastavam, TE-C390-siddhantam, TE-C390-sutram, TE-C390-prakriya, TE-C390-nirvacincu

\chapter{Facts and Theories}
\label{ch:te-facts-and-theories}
\hlchaptermodality{390}

\begin{chapteropening}
\textbf{By the end of this chapter:} \emph{I can talk about facts, theories, principles and processes.}

Complete the last lesson of chapter 390: I can talk about facts, theories, principles and processes.
\end{chapteropening}

\section[\texorpdfstring{\te{వాస్తవం} (vāstavaṁ)}{వాస్తవం (vāstavaṁ)}]{\te{వాస్తవం} (vāstavaṁ) — a fact; reality}

\label{lesson:TE-C390-vastavam}



\begin{quote}
Before the new one: say the Telugu for a likeness, then the Telugu for a distinction.
\end{quote}

\subsection*{You'll want to know: \te{వాస్తవం}}
\textbf{\te{వాస్తవం}} — \emph{vāstavaṁ} — \textquotedblleft{}a fact; reality\textquotedblright{}.

\begin{quote}
\te{అది} \te{వాస్తవం} \te{కాదు}, \te{ఒక} \te{కథ} \te{మాత్రమే}.
\end{quote}

\begin{quote}
\emph{adi vāstavaṁ kādu, oka katha mātramē.}
\end{quote}

\begin{quote}
That is not a fact, only a story.
\end{quote}

Plural: \te{వాస్తవాలు}, \emph{vāstavālu}.

\te{వాస్తవానికి}, in fact, as a matter of fact.

\begin{cousinweb}
\textbf{\te{వాస్తవం}} (\emph{vāstavaṁ}) is borrowed from Sanskrit \emph{vāstava}, \textquotedblleft{}substantial, real, true\textquotedblright{} (Monier-Williams), from \emph{vastu}, \textquotedblleft{}a thing\textquotedblright{}.
\end{cousinweb}

\begin{grammarlens}[title={What you've built}]
One more word for giving reasons and opinions.
\end{grammarlens}

\subsection*{Your turn}
\begin{itemize}
\raggedright
  \item \emph{Say it:} \emph{vāstavaṁ}
  \item \emph{Say it:} \emph{vāstavaṁ}, once more
  \item \emph{Say it:} \emph{adi vāstavaṁ kādu, oka katha mātramē.} — That is not a fact, only a story.
  \item \emph{Recall:} say the Telugu for a likeness, then the Telugu for a distinction, then say \emph{vāstavaṁ} again
\end{itemize}

\begin{tcolorbox}[breakable,colback=teal!4,colframe=teal!35!black,title={Before you move on}]
What does \te{వాస్తవం} mean? (\textquotedblleft{}A fact; reality\textquotedblright{}.) Say it once more.
\end{tcolorbox}
\section[\texorpdfstring{\te{సిద్ధాంతం} (siddhāntaṁ)}{సిద్ధాంతం (siddhāntaṁ)}]{\te{సిద్ధాంతం} (siddhāntaṁ) — a theory; a doctrine}

\label{lesson:TE-C390-siddhantam}



\begin{quote}
Before the new one: say the Telugu for a distinction, then the Telugu for a fact.
\end{quote}

\subsection*{You'll want to know: \te{సిద్ధాంతం}}
\textbf{\te{సిద్ధాంతం}} — \emph{siddhāntaṁ} — \textquotedblleft{}a theory; a doctrine\textquotedblright{}.

\begin{quote}
\te{ఆ} \te{సిద్ధాంతాన్ని} \te{ఇంకా} \te{ఎవరూ} \te{నిరూపించలేదు}.
\end{quote}

\begin{quote}
\emph{ā siddhāntānni iṅkā evarū nirūpiñcalēdu.}
\end{quote}

\begin{quote}
No one has proved that theory yet.
\end{quote}

Plural: \te{సిద్ధాంతాలు}, \emph{siddhāntālu}.

Also a principle a person lives by.

\begin{cousinweb}
\textbf{\te{సిద్ధాంతం}} (\emph{siddhāntaṁ}) is borrowed from Sanskrit \emph{siddhānta}, \textquotedblleft{}established end, doctrine\textquotedblright{} (CDIAL 13404), from \emph{siddha}, \textquotedblleft{}established\textquotedblright{}, and \emph{anta}, \textquotedblleft{}end\textquotedblright{}.
\end{cousinweb}

\begin{grammarlens}[title={What you've built}]
One more word for giving reasons and opinions.
\end{grammarlens}

\subsection*{Your turn}
\begin{itemize}
\raggedright
  \item \emph{Say it:} \emph{siddhāntaṁ}
  \item \emph{Say it:} \emph{siddhāntaṁ}, once more
  \item \emph{Say it:} \emph{ā siddhāntānni iṅkā evarū nirūpiñcalēdu.} — No one has proved that theory yet.
  \item \emph{Recall:} say the Telugu for a distinction, then the Telugu for a fact, then say \emph{siddhāntaṁ} again
\end{itemize}

\begin{tcolorbox}[breakable,colback=teal!4,colframe=teal!35!black,title={Before you move on}]
What does \te{సిద్ధాంతం} mean? (\textquotedblleft{}A theory; a doctrine\textquotedblright{}.) Say it once more.
\end{tcolorbox}
\section[\texorpdfstring{\te{సూత్రం} (sūtraṁ)}{సూత్రం (sūtraṁ)}]{\te{సూత్రం} (sūtraṁ) — a principle; a formula}

\label{lesson:TE-C390-sutram}



\begin{quote}
Before the new one: say the Telugu for a fact, then the Telugu for a theory.
\end{quote}

\subsection*{You'll want to know: \te{సూత్రం}}
\textbf{\te{సూత్రం}} — \emph{sūtraṁ} — \textquotedblleft{}a principle; a formula\textquotedblright{}.

\begin{quote}
\te{గణితంలో} \te{ప్రతి} \te{సూత్రానికి} \te{ఒక} \te{కారణం} \te{ఉంటుంది}.
\end{quote}

\begin{quote}
\emph{gaṇitaṁlō prati sūtrāniki oka kāraṇaṁ uṇṭundi.}
\end{quote}

\begin{quote}
Every formula in mathematics has a reason behind it.
\end{quote}

Plural: \te{సూత్రాలు}, \emph{sūtrālu}.

First a thread, then a short rule that holds a lot together.

\begin{cousinweb}
\textbf{\te{సూత్రం}} (\emph{sūtraṁ}) is borrowed from Sanskrit \emph{sūtra}, \textquotedblleft{}a thread, a cord\textquotedblright{} (CDIAL 13561), which Sanskrit also used for a short rule or formula (Monier-Williams).
\end{cousinweb}

\begin{grammarlens}[title={What you've built}]
One more word for giving reasons and opinions.
\end{grammarlens}

\subsection*{Your turn}
\begin{itemize}
\raggedright
  \item \emph{Say it:} \emph{sūtraṁ}
  \item \emph{Say it:} \emph{sūtraṁ}, once more
  \item \emph{Say it:} \emph{gaṇitaṁlō prati sūtrāniki oka kāraṇaṁ uṇṭundi.} — Every formula in mathematics has a reason behind it.
  \item \emph{Recall:} say the Telugu for a fact, then the Telugu for a theory, then say \emph{sūtraṁ} again
\end{itemize}

\begin{tcolorbox}[breakable,colback=teal!4,colframe=teal!35!black,title={Before you move on}]
What does \te{సూత్రం} mean? (\textquotedblleft{}A principle; a formula\textquotedblright{}.) Say it once more.
\end{tcolorbox}
\section[\texorpdfstring{\te{ప్రక్రియ} (prakriya)}{ప్రక్రియ (prakriya)}]{\te{ప్రక్రియ} (prakriya) — a process; a procedure}

\label{lesson:TE-C390-prakriya}



\begin{quote}
Before the new one: say the Telugu for a theory, then the Telugu for a principle.
\end{quote}

\subsection*{You'll want to know: \te{ప్రక్రియ}}
\textbf{\te{ప్రక్రియ}} — \emph{prakriya} — \textquotedblleft{}a process; a procedure\textquotedblright{}.

\begin{quote}
\te{ఆ} \te{ప్రక్రియకు} \te{రెండు} \te{నెలలు} \te{పట్టింది}.
\end{quote}

\begin{quote}
\emph{ā prakriyaku reṇḍu nelalu paṭṭindi.}
\end{quote}

\begin{quote}
That process took two months.
\end{quote}

Any step-by-step official or technical process. \te{పట్టు} here means to take time.

\begin{cousinweb}
\textbf{\te{ప్రక్రియ}} (\emph{prakriya}) is borrowed from Sanskrit \emph{prakriyā}, \textquotedblleft{}production; procedure, way, manner\textquotedblright{} (Monier-Williams); Telugu shortens the final long a.
\end{cousinweb}

\begin{grammarlens}[title={What you've built}]
One more word for giving reasons and opinions.
\end{grammarlens}

\subsection*{Your turn}
\begin{itemize}
\raggedright
  \item \emph{Say it:} \emph{prakriya}
  \item \emph{Say it:} \emph{prakriya}, once more
  \item \emph{Say it:} \emph{ā prakriyaku reṇḍu nelalu paṭṭindi.} — That process took two months.
  \item \emph{Recall:} say the Telugu for a theory, then the Telugu for a principle, then say \emph{prakriya} again
\end{itemize}

\begin{tcolorbox}[breakable,colback=teal!4,colframe=teal!35!black,title={Before you move on}]
What does \te{ప్రక్రియ} mean? (\textquotedblleft{}A process; a procedure\textquotedblright{}.) Say it once more.
\end{tcolorbox}
\section[\texorpdfstring{\te{నిర్వచించు} (nirvaciñcu)}{నిర్వచించు (nirvaciñcu)}]{\te{నిర్వచించు} (nirvaciñcu) — to define}

\label{lesson:TE-C390-nirvacincu}



\begin{quote}
Before the new one: say the Telugu for a principle, then the Telugu for a process.
\end{quote}

\subsection*{You'll want to know: \te{నిర్వచించు}}
\textbf{\te{నిర్వచించు}} — \emph{nirvaciñcu} — \textquotedblleft{}to define\textquotedblright{}.

\begin{quote}
\te{ఉపాధ్యాయుడు} \te{ఆ} \te{మాటను} \te{స్పష్టంగా} \te{నిర్వచించాడు}.
\end{quote}

\begin{quote}
\emph{upādhyāyuḍu ā māṭanu spaṣṭaṅgā nirvaciñcāḍu.}
\end{quote}

\begin{quote}
The teacher defined that word clearly.
\end{quote}

Past (I …): \te{నిర్వచించాను}, \emph{nirvaciñcānu}.

The noun is \te{నిర్వచనం}, a definition.

\begin{cousinweb}
\textbf{\te{నిర్వచించు}} (\emph{nirvaciñcu}) is built on Sanskrit \emph{nirvacana}, \textquotedblleft{}speaking out, pronouncing\textquotedblright{} (Monier-Williams), with the Telugu verb ending -\te{ఇంచు}.
\end{cousinweb}

\begin{grammarlens}[title={What you've built}]
One more word for giving reasons and opinions.
\end{grammarlens}

\subsection*{Your turn}
\begin{itemize}
\raggedright
  \item \emph{Say it:} \emph{nirvaciñcu}
  \item \emph{Say it:} \emph{nirvaciñcu}, once more
  \item \emph{Say it:} \emph{upādhyāyuḍu ā māṭanu spaṣṭaṅgā nirvaciñcāḍu.} — The teacher defined that word clearly.
  \item \emph{Recall:} say the Telugu for a principle, then the Telugu for a process, then say \emph{nirvaciñcu} again
\end{itemize}

\begin{tcolorbox}[breakable,colback=teal!4,colframe=teal!35!black,title={Before you move on}]
What does \te{నిర్వచించు} mean? (\textquotedblleft{}To define\textquotedblright{}.) Say it once more.
\end{tcolorbox}
